% Einheit 1 von 3 – Potenzen (bank/potenzen-wurzeln/e1.jsonl, zusammenbau v0.1)
%% TODO Zweigzeile Teil 1: was hier gelernt wird (Satz in Schülersprache aus dem Einheitstitel) – Einheit 1
\einheitenkopf[e1]{Einheit 1 von 3 · Potenzen}
\zweigzeile{ab Kl. 5, je nach Buch bis Kl. 9 · P10 oft}
\setcounter{aufgabe}{11}

%% TODO Titel: Ich-kann-Satz fehlt in der Bank (Platzhalter: Kettenname) – Nr. 12
%% TODO Anweisung: ein Satz über der ersten Teilaufgabe fehlt in der Bank – Nr. 12
\begin{aufgabe}{Plus oder minus?}
\begin{teile}
\teil $(-5)^4$ – plus oder minus?\\ \kreuz{positiv} \\ \kreuz{negativ}
\end{teile}
\end{aufgabe}

%% TODO Titel: Ich-kann-Satz fehlt in der Bank (Platzhalter: Kettenname) – Nr. 13
%% TODO Anweisung: ein Satz über der ersten Teilaufgabe fehlt in der Bank – Nr. 13
\begin{aufgabe}{Bruch oder minus?}
\begin{teile}
\teil $7^{-2}$ – Bruch oder negative Zahl?\\ \kreuz{Bruch} \\ \kreuz{negative Zahl}
\end{teile}
\end{aufgabe}

%% TODO Titel: Ich-kann-Satz fehlt in der Bank (Platzhalter: Kettenname) – Nr. 14
%% TODO Anweisung: ein Satz über der ersten Teilaufgabe fehlt in der Bank – Nr. 14
\begin{aufgabe}{Potenzen}
%% TODO \rechenplatz{n}: Schrittzahl des Grundfalls fehlt in der Bank (nur bei fester Schreibform, 2.3 b)
\begin{teile}
\teil $7^4$ – hoch oder mal?\\ \kreuz{hoch} \\ \kreuz{mal}
\teil $2^6$ als Malkette – wie viel? \leerfeld
\teil $7^2$ und $7 \cdot 2$ – wie viel jeweils? \leerfeld und \leerfeld
\teil $14^2$ – wie viel? \leerfeld
\teil $23^2$ – wie viel? Rechne mit dem Taschenrechner. \leerfeld
\teil $(-4)^3$ – wie viel? \leerfeld
\teil $-8^2$ – wie viel? \leerfeld
\teil $0{,}6^2$ – wie viel? \leerfeld
\teil $2^7$ oder $5^3$ – welche Zahl ist größer?
\teil $5^3$, $2^6$, $9^2$ – welche Zahl ist die größte? Unterstreiche sie.
\teil $5^x = 625$ – welche Zahl ist $x$? x = \leerfeld
\teil $4^x = 4\,096$ – welche Zahl ist $x$? x = \leerfeld
\teil $10^x = 10\,000$ – welche Zahl ist $x$? x = \leerfeld
\teil $4^{-2}$ – als Bruch und als Dezimalzahl? \leerfeld = \leerfeld
\teil $4^{-3} \;\square\; 0{,}02$ – $<$, $=$ oder $>$?
\teil $2^x = 512$ – welche Zahl ist $x$? \hfill (P10 2020 OS) \\ x = \leerfeld
\end{teile}
\end{aufgabe}

%% TODO Titel: Ich-kann-Satz fehlt in der Bank (Platzhalter: Kettenname) – Nr. 15
%% TODO Anweisung: ein Satz über der ersten Teilaufgabe fehlt in der Bank – Nr. 15
\begin{aufgabe}{Potenzen von Brüchen}
\begin{teile}
\teil $\left(\frac{2}{5}\right)^2$ – wie viel? \leerfeld
\end{teile}
\end{aufgabe}

%% TODO Titel: Ich-kann-Satz fehlt in der Bank (Platzhalter: Kettenname) – Nr. 16
%% TODO Anweisung: ein Satz über der ersten Teilaufgabe fehlt in der Bank – Nr. 16
\begin{aufgabe}{Potenz gegen Dezimalzahl und Prozent vergleichen}
\begin{teile}
\teil $0{,}5^2 \;\square\; 20\,\%$ – $<$, $=$ oder $>$?
\end{teile}
\end{aufgabe}

%% TODO Titel: Ich-kann-Satz fehlt in der Bank (Platzhalter: Kettenname) – Nr. 17
%% TODO Anweisung: ein Satz über der ersten Teilaufgabe fehlt in der Bank – Nr. 17
\begin{aufgabe}{hoch null}
\begin{teile}
\teil $8^0$ – wie viel? \leerfeld
\end{teile}
\end{aufgabe}

%% TODO Titel: Ich-kann-Satz fehlt in der Bank (Platzhalter: Kettenname) – Nr. 18
%% TODO Anweisung: ein Satz über der ersten Teilaufgabe fehlt in der Bank – Nr. 18
\begin{aufgabe}{Potenzen – Fehler finden, Begründen}
\begin{teile}
\teil Mia rechnet $6^3 = 18$. Wo ist der Fehler?
\teil Begründe, warum $3^5$ größer ist als $5^3$.
\end{teile}
\end{aufgabe}
